\documentclass[../../main.tex]{subfiles}% 注意这里的文件路径不能用 ./main.tex ,否则用latexmk编译子文件会报错
\graphicspath{{\subfix{./image/}}{\subfix{../../image/}}} % 指定图片引用路径,后续可以直接使用图片文件名
% 如果图片存放在上级目录,则使用 ../../image/ 作为备用路径

% 例如:
% \begin{figure}[H]
% \centering
% \includegraphics[width=0.3\textwidth]{图.png}
% \caption{}
% \label{figure:图}
% \end{figure}
% 注意:上述\label{}一定要放在\caption{}之后,否则引用图片序号会只会显示??.

\begin{document}

\section{草稿二}


\begin{theorem}[扩张定理]\label{theorem:ch05-扩张定理}
设 $1\leqslant p\leqslant\infty$，$U\subseteq \mathbb{R}^n$是有界开集且 $\partial U$ 为 $C^{1}$。选取有界开集 $V$ 使 $U\subset \subset V$。则存在有界线性算子
\begin{align}
E:W^{1,p}(U)\to W^{1,p}(\mathbb{R}^{n})
\end{align}
使对每个 $u\in W^{1,p}(U)$：
\begin{enumerate}[(i)]
\item $Eu=u$ a.e. $x\in U$；
\item $\text{spt}\,Eu\subset \subset V$；
\item $\|Eu\|_{W^{1,p}(\mathbb{R}^{n})}\leqslant C\|u\|_{W^{1,p}(U)}$，其中常数 $C$ 只依赖于 $n,p,U,V$。
\end{enumerate}
我们称 $Eu$ 为 $u$ 到 $\mathbb{R}^{n}$ 的\textbf{扩张}。
\end{theorem}
\begin{proof}

\noindent$\mathbf{Step}\bm{1}$.{\heiti 平坦边界处的半空间与反射公式.}

固定 \(x_0 \in \partial U\). 当 \(x_0 = 0\)，且\(\partial U\) 在 \(x_0\) 附近平坦，位于平面 \(\{x_n=0\}\) 内时，即存在 \(r>0\)，使得开球 \(B \triangleq B(0,r)\) 满足
\begin{align*}
\begin{cases}
B^{+}\triangleq B\cap\{x_{n}> 0\}\subset\overline{U},\\
B^{-}\triangleq B\cap\{x_{n}< 0\}\subset\mathbb{R}^{n}\backslash U.
\end{cases}
\end{align*}

% \begin{figure}[H]
% \centering
% \includegraphics[width=0.55\textwidth]{ch05-half-ball.png}
% \caption{边界处的半球}
% \label{figure:ch05-边界半球}
% \end{figure}

暂时再设 $u\in C^{1}(\overline{U})$。于是定义
\begin{align}
\overline{u}(x)\triangleq
\begin{cases}
u\left( x \right) ,&		x\in B^+,\\
-3u\left( x_1,\cdots ,x_{n-1},-x_n \right) +4u\left( x_1,\cdots ,x_{n-1},-\frac{x_n}{2} \right) ,&		x\in B^-.\\
\end{cases}
\label{align:ch05-反射延拓}
\end{align}
这称为把 $u$ 从 $B^{+}$ \textbf{高阶反射}到 $B^{-}$。

\noindent$\mathbf{Step}\bm{2}$.{\heiti 反射延拓的 \(C^1\) 正则性.}

我们断言
\begin{align*}
\overline{u}\in C^{1}(B).
\end{align*}
首先，\(\overline{u}\) 在 \(B^+\) 上等于 \(u\)，在 \(B^-\) 上由 \eqref{align:ch05-反射延拓}式定义。由于 \(u \in C^1(\overline{U})\)，且当 \(x \in B^-\) 时，点 \((x',-x_n)\) 和 \(\left(x',-\frac{x_n}{2}\right)\) 均属于 \(B^+ \subseteq \overline{U}\)，所以 \(\overline{u}\) 在 \(B^+\) 内部和 \(B^-\) 内部均属于 \(C^1\)。因此只需验证 \(\overline{u}\) 及其一阶偏导数在边界 \(\{x_n=0\}\) 处连续。

记$B^{0}\triangleq B\cap\{x_n=0\}$,
固定
\[
x_0=(x_0',0)\in B^{0}.
\]
由于 \(u\in C^{1}(\overline U)\)，\(u\) 在紧集 \(\overline{B^{+}}\cup B^{0}\) 上一致连续。因此当
\[
(x',x_n)\to(x_0',0)
\]
时，有联合极限
\[
u(x',x_n)\to u(x_0',0),\qquad
u(x',-x_n)\to u(x_0',0),\qquad
u\left(x',-\frac{x_n}{2}\right)\to u(x_0',0).
\]
于是
\[
\overline u(x)
=-3u(x',-x_n)+4u\left(x',-\frac{x_n}{2}\right)
\to -3u(x_0',0)+4u(x_0',0)
=u(x_0',0)
=\overline u(x_0).
\]
故 \(\overline u\) 在 \(B\) 上连续。

其次计算一阶偏导数。在 \(B^{+}\) 上，
\[
D_i\overline u(x)=D_i u(x),\qquad i=1,\cdots,n.
\]
在 \(B^{-}\) 上，直接求导得
\[
D_i\overline u(x)=
\begin{cases}
-3D_i u(x',-x_n)+4D_i u\left(x',-\dfrac{x_n}{2}\right), & i=1,\cdots,n-1,\\[6pt]
3D_n u(x',-x_n)-2D_n u\left(x',-\dfrac{x_n}{2}\right), & i=n.
\end{cases}
\]
定义 \(g_i:B\to\mathbb{R}\) 为
\[
g_i(x)\triangleq
\begin{cases}
D_i u(x), & x\in B^{+}\cup B^{0},\\[4pt]
-3D_i u(x',-x_n)+4D_i u\left(x',-\dfrac{x_n}{2}\right), & x\in B^{-},\ i<n,\\[6pt]
3D_n u(x',-x_n)-2D_n u\left(x',-\dfrac{x_n}{2}\right), & x\in B^{-},\ i=n.
\end{cases}
\]
则 \(g_i\) 在 \(B\setminus B^{0}\) 上连续，且
\[
D_i\overline u(x)=g_i(x),\qquad x\in B^{+}\cup B^{-}.
\]

下面证明 \(g_i\) 在 \(B^{0}\) 上也连续。任取 \(x_0=(x_0',0)\in B^{0}\)。
若 \(x\in B^{+}\)，则
\[
g_i(x)=D_i u(x).
\]
由于 \(D_i u\) 在 \(\overline{B^{+}}\) 上一致连续，故
\[
g_i(x)\to D_i u(x_0',0)=g_i(x_0).
\]

若 \(x\in B^{-}\)，则当 \(x\to x_0\) 时，
\[
(x',-x_n)\to(x_0',0),\qquad
\left(x',-\frac{x_n}{2}\right)\to(x_0',0).
\]
由 \(D_i u\) 在 \(\overline{B^{+}}\) 上的一致连续性，
\[
D_i u(x',-x_n)\to D_i u(x_0',0),\qquad
D_i u\left(x',-\frac{x_n}{2}\right)\to D_i u(x_0',0).
\]
因此
\[
g_i(x)\to
\begin{cases}
-3D_i u(x_0',0)+4D_i u(x_0',0)=D_i u(x_0',0), & i<n,\\[4pt]
3D_n u(x_0',0)-2D_n u(x_0',0)=D_n u(x_0',0), & i=n.
\end{cases}
\]
故
\[
g_i(x)\to D_i u(x_0',0)=g_i(x_0).
\]
所以 \(g_i\) 在 \(B^{0}\) 上也连续，从而 \(g_i\) 在 \(B\) 上连续。

下面证明在边界点 \(x_0=(x_0',0)\in B^{0}\) 处，\(\overline u\) 的各偏导数存在且等于 \(g_i(x_0)\)。

对 \(i=1,\cdots,n-1\)，当 \(h\to0\) 时，\(x_0+he_i\in B^{0}\)，于是
\[
\frac{\overline u(x_0+he_i)-\overline u(x_0)}{h}
=
\frac{u(x_0'+he_i,0)-u(x_0',0)}{h}.
\]
由于 \(u(\cdot,0)\in C^{1}\)，所以
\[
\frac{\overline u(x_0+he_i)-\overline u(x_0)}{h}
\to D_i u(x_0',0)=g_i(x_0).
\]
因此
\[
D_i\overline u(x_0)=g_i(x_0),\qquad i=1,\cdots,n-1.
\]

对 \(i=n\)，先取 \(h>0\)。函数
\[
t\mapsto\overline u(x_0+te_n)
\]
在 \([0,h]\) 上连续，在 \((0,h)\) 内可微。由中值定理，存在 \(\theta_h\in(0,h)\)，使得
\[
\frac{\overline u(x_0+he_n)-\overline u(x_0)}{h}
=D_n\overline u(x_0+\theta_h e_n).
\]
当 \(h\to0^{+}\) 时，\(\theta_h\to0\)，且 \(x_0+\theta_h e_n\in B^{+}\)。由 \(g_n\) 在边界处的连续性，
\[
D_n\overline u(x_0+\theta_h e_n)
=g_n(x_0+\theta_h e_n)
\to g_n(x_0).
\]
再取 \(h<0\)。同理，在 \(B^{-}\) 上用中值定理，存在 \(\theta_h\in(h,0)\)，使得
\[
\frac{\overline u(x_0+he_n)-\overline u(x_0)}{h}
=D_n\overline u(x_0+\theta_h e_n)
=g_n(x_0+\theta_h e_n)
\to g_n(x_0).
\]
因此两个单侧差商极限相同，故
\[
D_n\overline u(x_0)=g_n(x_0).
\]
于是
\[
D_i\overline u(x)=g_i(x),\qquad x\in B,\quad i=1,\cdots,n.
\]
由于每个 \(D_i\overline u=g_i\) 都在 \(B\) 上连续，\(\overline u\) 的所有一阶偏导数在 \(B\) 上存在且连续。由多元函数可微的充分条件，\(\overline u\) 在 \(B\) 上可微。因此
\[
\overline u\in C^{1}(B).
\]

\noindent$\mathbf{Step}\bm{3}$.{\heiti 半空间反射的 \(W^{1,p}\) 估计.}

现在证明
\begin{align*}
\|\overline{u}\|_{W^{1,p}(B)}\leqslant C\|u\|_{W^{1,p}(B^{+})},
\end{align*}
其中常数 $C$ 不依赖于 $u$。
当$1\leqslant p<\infty$时,先估计 \(L^p\) 部分. 在 \(B^+\) 上 \(\overline{u}=u\), 故
\begin{align*}
\int_{B^+}|\overline{u}|^p\,\mathrm{d}x=\int_{B^+}|u|^p\,\mathrm{d}x.
\end{align*}
在 \(B^-\) 上, 由 \eqref{align:ch05-反射延拓}式和\hyperref[Basis of Analytics-proposition:凸性不等式]{凸性不等式}得
\begin{align*}
|\overline{u}(x_1,\cdots,x_n)|^p
\leqslant C_p\left(3^p|u(x_1,\cdots,x_{n-1},-x_n)|^p
+4^p\left|u\left(x_1,\cdots,x_{n-1},-\frac{x_n}{2}\right)\right|^p\right),
\end{align*}
其中$C_p\triangleq 2^{p-1}$.因此
\begin{align}\label{eq::jiojiojiojioqwio2i3333-1}
\int_{B^-}|\overline{u}|^p\,\mathrm{d}x
\leqslant C_p\left(3^p\int_{B^-}|u(x_1,\cdots,x_{n-1},-x_n)|^p\,\mathrm{d}x
+4^p\int_{B^-}\left|u\left(x_1,\cdots,x_{n-1},-\frac{x_n}{2}\right)\right|^p\,\mathrm{d}x\right).
\end{align}
对第一个积分作变量替换 \(x_n\mapsto -x_n\),其他不变. 因为 \(B\) 关于 \(\{x_n=0\}\) 对称, \(B^-\) 变为 \(B^+\), 且Jacobi行列式绝对值为 \(1\), 所以
\begin{align}\label{eq::jiojiojiojioqwio2i3333-0}
\int_{B^-}|u(x_1,\cdots,x_{n-1},-x_n)|^p\,\mathrm{d}x
=\int_{B^+}|u(x_1,\cdots,x_n)|^p\,\mathrm{d}x.
\end{align}
对第二个积分作变量替换 \(x_n\mapsto -\frac{x_n}{2}\),其他不变. 此时 \(B^-\) 变为
\begin{align*}
\{x_1^2+\cdots+x_{n-1}^2+4x_n^2<r^2,\ x_n>0\},
\end{align*}
它包含于 \(B^+\), 且该变量替换的Jacobi行列式绝对值为 \(2\), 因此
\begin{align}\label{eq::jiojiojiojioqwio2i3333-2}
\int_{B^-}\left|u\left(x_1,\cdots,x_{n-1},-\frac{x_n}{2}\right)\right|^p\,\mathrm{d}x
=2\int_{\{x_1^2+\cdots+x_{n-1}^2+4x_n^2<r^2,\ x_n>0\}}|u(x_1,\cdots,x_n)|^p\,\mathrm{d}x
\leqslant 2\int_{B^+}|u|^p\,\mathrm{d}x.
\end{align}
合并\eqref{eq::jiojiojiojioqwio2i3333-1}\eqref{eq::jiojiojiojioqwio2i3333-0}\eqref{eq::jiojiojiojioqwio2i3333-2}式得
\begin{align*}
\int_{B^-}{\left| \overline{u} \right|^p\mathrm{d}x}\leqslant C_p\left( 3^p+2\cdot 4^p \right) \int_{B^+}{\left| u \right|^p\mathrm{d}x}.
\end{align*}
从而
\begin{align}
\left\| \overline{u} \right\| _{L^p\left( B \right)}^{p}&=\int_B{\left| \overline{u} \right|^p\mathrm{d}x}=\int_{B^+}{\left| u \right|^p\mathrm{d}x}+\int_{B^-}{\left| \overline{u} \right|^p\mathrm{d}x}
\nonumber \\
&\leqslant \int_{B^+}{\left| u \right|^p\mathrm{d}x}+C_p\left( 3^p+2\cdot 4^p \right) \int_{B^+}{\left| u \right|^p\mathrm{d}x}
\nonumber \\
&=\left[ 1+C_p\left( 3^p+2\cdot 4^p \right) \right] \left\| u \right\| _{L^p\left( B^+ \right)}^{p}.\label{eq:::HJiuhj3u892HFFHJIOEJO}
\end{align}
再估计一阶导数. 对 \(i=1,\cdots,n-1\), 在 \(B^-\) 上
\begin{align*}
D_i\overline{u}(x_1,\cdots,x_n)
=-3D_i u(x_1,\cdots,x_{n-1},-x_n)
+4D_i u\left(x_1,\cdots,x_{n-1},-\frac{x_n}{2}\right),
\end{align*}
利用与\eqref{eq::jiojiojiojioqwio2i3333-1}\eqref{eq::jiojiojiojioqwio2i3333-0}式同样的变量替换给出
\begin{align*}
\int_{B^-}|D_i\overline{u}|^p\,\mathrm{d}x
\leqslant C_p(3^p+2\cdot 4^p)\int_{B^+}{\left| D_iu \right|^p\mathrm{d}x}.
\end{align*}
对 \(i=n\), 在 \(B^-\) 上
\begin{align*}
D_n\overline{u}(x_1,\cdots,x_n)
=3D_n u(x_1,\cdots,x_{n-1},-x_n)
-2D_n u\left(x_1,\cdots,x_{n-1},-\frac{x_n}{2}\right),
\end{align*}
再利用与\eqref{eq::jiojiojiojioqwio2i3333-1}\eqref{eq::jiojiojiojioqwio2i3333-0}式同样的变量替换给出
\begin{align*}
\int_{B^-}|D_n\overline{u}|^p\,\mathrm{d}x
\leqslant C_p(3^p+2\cdot 2^p)\int_{B^+}{\left| D_nu \right|^p\mathrm{d}x}.
\end{align*}
在 \(B^+\) 上 \(\overline{u}=u\), 所以对每个 \(i=1,\cdots,n\),有
\begin{align}
\left\| D_i\overline{u} \right\| _{L^p\left( B \right)}^{p}&=\int_B{\left| D_i\overline{u} \right|^p\mathrm{d}x}=\int_{B^+}{\left| D_iu \right|^p\mathrm{d}x}+\int_{B^-}{\left| D_i\overline{u} \right|^p\mathrm{d}x}
\nonumber \\
&\leqslant \int_{B^+}{\left| D_iu \right|^p\mathrm{d}x}+C_p\left( 3^p+2\cdot 4^p \right) \int_{B^+}{\left| D_iu \right|^p\mathrm{d}x}
\nonumber \\
&=\left[ 1+C_p\left( 3^p+2\cdot 4^p \right) \right] \left\| D_iu \right\| _{L^p\left( B^+ \right)}^{p}.\label{eq:::HJiuhj3u892HFFHJIOEJO-1}
\end{align}
取$C\triangleq \left[ 1+C_p(3^p+2\cdot 4^p) \right] ^{\frac{1}{p}}$,则由\eqref{eq:::HJiuhj3u892HFFHJIOEJO}\eqref{eq:::HJiuhj3u892HFFHJIOEJO-1}式可得
\begin{align*}
\|\overline{u}\|_{W^{1,p}(B)}^p
=\|\overline{u}\|_{L^p(B)}^p+\sum\limits_{i=1}^n\|D_i\overline{u}\|_{L^p(B)}^p\leqslant C^p\left(\|u\|_{L^p(B^+)}^p+\sum\limits_{i=1}^n\|D_i u\|_{L^p(B^+)}^p\right)=C^p\|u\|_{W^{1,p}(B^+)}^p.
\end{align*}
开 \(p\) 次方得
\begin{align*}
\|\overline{u}\|_{W^{1,p}(B)}
\leqslant C\|u\|_{W^{1,p}(B^+)}.
\end{align*}
若 \(p=\infty\), 直接取本性上界得
\begin{gather*}
\left| \overline{u} \right|\leqslant 3\left| u \right|+4\left| u \right|\leqslant 7\left\| u \right\| _{L^{\infty}\left( B^+ \right)},
\\
\left| D_i\overline{u} \right|\leqslant 3\left| D_iu \right|+4\left| D_iu \right|\leqslant 7\left\| D_iu \right\| _{L^{\infty}\left( B^+ \right)},\quad i=1,2,\cdots ,n.
\end{gather*}
取$C=7$,则有
\begin{align*}
\|\overline{u}\|_{W^{1,\infty}(B)}
\leqslant C\|u\|_{W^{1,\infty}(B^+)}.
\end{align*}

% \begin{figure}[H]
% \centering
% \includegraphics[width=0.8\textwidth]{8e1ffee58f4e42d8a67b7ea75f26d99d.png}
% \caption{拉直边界}
% \label{figure:ch05-拉直边界}
% \end{figure}

\noindent$\mathbf{Step}\bm{4}$.{\heiti 一般边界点的拉直.}

固定 \(x_0 \in \partial U\). 设 \(\nu(x_0)\) 为 \(x_0\) 处关于 \(U\) 的单位外法向量. 取$n$阶正交矩阵 \(Q \), 使得 \(Q^T\) 的最后一列为 \(-\nu(x_0)\). 定义平移与旋转变换
\begin{align}
z &\triangleq Q(x - x_0). \label{eq:ch05-straighten-coord}
\end{align}
在此新坐标系 \(z\) 下, 点 \(x_0\) 对应 \(z = 0\), 且 \(U\) 在 \(x_0\) 处的内法向量指向 \(z_n\) 轴的正方向. 由\cref{definition:appC-Ck边界}, 存在 \(r > 0\) 和 \(C^1\) 函数 \(\gamma: \mathbb{R}^{n-1} \to \mathbb{R}\), 使得
\begin{align}
U \cap B_z(0, r) &= \{z \in B_z(0, r) \mid z_n > \gamma(z_1, \cdots, z_{n-1})\}. \label{eq:ch05-boundary-def}
\end{align}
这里 \(z' = (z_1, \cdots, z_{n-1})\). 由于 \(z_n\) 轴沿外法向, 故 \(\gamma(0) = 0\) 且 \(\nabla \gamma(0) = 0\).

定义拉直映射 \(\Phi_1: B_z(0, r) \to \mathbb{R}^n\) 为
\begin{align*}
\Phi_1(z) &\triangleq (z_1, \cdots, z_{n-1}, z_n - \gamma(z_1, \cdots, z_{n-1})).
\end{align*}
其逆映射 \(\Psi_1: \Phi_1(B_z(0, r)) \to \mathbb{R}^n\) 为
\begin{align*}
\Psi_1(y) &\triangleq (y_1, \cdots, y_{n-1}, y_n + \gamma(y_1, \cdots, y_{n-1})).
\end{align*}
直接计算可得 \(\Phi_1(\Psi_1(y)) = y\), 且 \(\Psi_1(\Phi_1(z)) = z\). 因此 \(\Phi_1\) 与 \(\Psi_1\) 互为逆映射. 由于 \(\gamma \in C^1\), 故 \(\Phi_1, \Psi_1 \in C^1\), 它们均为 \(C^1\) 微分同胚，其Jacobi矩阵分别为
\begin{align*}
D\Phi _1\left( z \right) =\left( \begin{matrix}
I_{n-1}&		0\\
-\nabla \gamma \left( z' \right) ^T&		1\\
\end{matrix} \right) ,\quad D\Psi _1\left( z \right) =\left( \begin{matrix}
I_{n-1}&		0\\
\nabla \gamma \left( z' \right) ^T&		1\\
\end{matrix} \right) .
\end{align*}
故 \(\left| \det D\Phi _1(z) \right|=\left| \det D\Psi _1(y) \right|=1\).

现在先在 \(y\)坐标系下选取开球 \(B\). 令 \(y_0 \triangleq \Phi_1(0)=(0,\cdots,0)\). 
因为 \(\Psi_1\) 连续且 \(\Psi_1(0)=0\)，所以对任意 \(\varepsilon>0\)，存在 \(\rho_\varepsilon>0\)，使得
\begin{align*}
y\in B(0,\rho_\varepsilon)\Longrightarrow \|\Psi_1(y)\|<\varepsilon.
\end{align*}
又因为 \(Q\) 是正交矩阵，所以\(Q^T\) 保持欧氏范数，从而
\begin{align*}
W_\varepsilon\triangleq x_0+Q^T\Psi_1(B(0,\rho_\varepsilon))
\subseteq x_0+Q^TB(0,\varepsilon)
=B(x_0,\varepsilon).
\end{align*}
由于 \(x_0\in\partial U\subseteq \overline U\subseteq V\)，而 \(V\) 是开集，故存在 \(\delta>0\) 使得
\begin{align*}
B(x_0,\delta)\subseteq V.
\end{align*}
取$\varepsilon_1=\min\left\{r,\frac{\delta}{2}\right\}.$
则存在 \(\rho\triangleq \rho_{\varepsilon_1}>0\)，使得
\begin{align*}
W\triangleq W_{\varepsilon_1}\subseteq B(x_0,\varepsilon_1)\subseteq B(x_0,r)\cap B\left(x_0,\frac{\delta}{2}\right)\subseteq B(x_0,\delta)\subseteq V.
\end{align*}
故开球$B \triangleq B(y_0, \rho)$的原像
\begin{align*}
W\triangleq x_0+Q^T\Psi _1\left( B \right)
\end{align*}
包含在$B\left( x_0,r \right)$中,且$W\subset \subset V$.然后定义原坐标下的拉直映射 \(\Phi: W \to B\) 为
\begin{align*}
\Phi(x) &\triangleq \Phi_1(Q(x - x_0)),
\end{align*}
其逆映射 \(\Psi: B \to W\) 为
\begin{align*}
\Psi(y) &\triangleq x_0 + Q^T \Psi_1(y).
\end{align*}
由定义可知 \(\Phi\) 与 \(\Psi\) 互为逆映射, 且均为 \(C^1\) 微分同胚. 并且由\hyperref[theorem:向量求导基本性质]{向量求导基本性质\ref{theorem:向量求导基本性质-4}}知
\begin{gather}
\Phi(W) = B, \quad \Psi(B) = W,\label{eq::jifjiojiojiojiomfijsfiojjsifjioejiojmwfiejfiwmfiowijijwfio-1}\\
\left| \det D\Phi \left( x \right) \right|=\left| \det \left( D\Phi _1\left( Q\left( x-x_0 \right) \right) Q \right) \right|=\left| \det D\Phi _1\left( Q\left( x-x_0 \right) \right) \right|\cdot \left| \det Q \right|=1,\label{eq::jifjiojiojiojiomfijsfiojjsifjioejiojmwfiejfiwmfiowijijwfio-2}
\\
\left| \det D\Psi \left( y \right) \right|=\left| \det \left( Q^TD\Psi _1\left( y \right) \right) \right|=\left| \det Q^T \right|\cdot \left| \det D\Psi _1\left( y \right) \right|=1.\label{eq::jifjiojiojiojiomfijsfiojjsifjioejiojmwfiejfiwmfiowijijwfio-3}
\end{gather}

下面验证 \(\Phi\) 将边界拉直. 若 \(x \in W \cap U\), 则 \(z = Q(x - x_0) \in B_z(0, r)\), 且 \(z_n > \gamma(z')\). 于是
\begin{align*}
\Phi_n(x) &= (\Phi_1(z))_n = z_n - \gamma(z') > 0,
\end{align*}
即 \(\Phi(x) \in B^+\). 因此 \(\Phi(W \cap U) \subseteq B^+\). 反之, 若 \(y \in B^+\), 则 \(y_n > 0\). 令 \(x = \Psi(y)\), 则 \(z = Q(x - x_0) = \Psi_1(y)\). 于是
\begin{align*}
z_n &= y_n + \gamma(y') > \gamma(y') = \gamma(z'),
\end{align*}
故 \(x \in W \cap U\). 因此 \(\Phi(W \cap U) = B^+\). 若 \(x \in W \cap \partial U\), 则 \(z_n = \gamma(z')\), 于是 \(\Phi_n(x) = 0\), 即 \(\Phi(x) \in B \cap \{y_n = 0\}\). 因此 \(\Phi(W \cap \partial U) = B \cap \{y_n = 0\}\).

\noindent$\mathbf{Step}\bm{5}$.{\heiti 拉直后的局部延拓与估计.}

下面在 \(B\) 中作反射延拓. 先设 \(u \in C^1(\overline{U})\). 定义
\begin{align*}
u'(y) \triangleq u(\Psi(y)), \quad y \in B^+.
\end{align*}
对 \(u'\) 应用$\mathbf{Step}\bm{1-3}$中的高阶反射延拓, 得到 \(\overline{u}' \in C^1(B)\), 满足
\begin{gather}
\overline{u}'(y) = u'(y), \quad y \in B^+, \label{eq:ch05-ext-in-ball-1} \\
\|\overline{u}'\|_{W^{1,p}(B)} \leqslant C \|u'\|_{W^{1,p}(B^+)}. \label{eq:ch05-ext-in-ball-2}
\end{gather}

定义 \(\overline{u}(x) \triangleq \overline{u}'(\Phi(x))\), \(x \in W\). 则 \(\overline{u} \in C^1(W)\). 若 \(x \in W \cap U\), 则 \(\Phi(x) \in B^+\), 于是
\begin{align*}
\overline{u}(x) &= \overline{u}'(\Phi(x)) = u'(\Phi(x)) = u(\Psi(\Phi(x))) = u(x).
\end{align*}
故 \(\overline{u}|_{W \cap U} = u\). 由变量替换 \(y = \Phi(x)\), \(x = \Psi(y)\) 及\eqref{eq::jifjiojiojiojiomfijsfiojjsifjioejiojmwfiejfiwmfiowijijwfio-1}\eqref{eq::jifjiojiojiojiomfijsfiojjsifjioejiojmwfiejfiwmfiowijijwfio-2}\eqref{eq::jifjiojiojiojiomfijsfiojjsifjioejiojmwfiejfiwmfiowijijwfio-3}式, 得
\begin{gather}
\int_W |\overline{u}(x)|^p \mathrm{d}x = \int_B |\overline{u}'(y)|^p |\det D\Psi(y)| \mathrm{d}y = \int_B |\overline{u}'(y)|^p \mathrm{d}y,\label{eq::jijioqwio93283r3r23}\\
\int_{B^+} |u'(y)|^p \mathrm{d}y = \int_{W \cap U} |u(x)|^p |\det D\Phi(x)| \mathrm{d}x = \int_{W \cap U} |u(x)|^p \mathrm{d}x.\label{eq::jijioqwio93283r3r23-111}
\end{gather}
由链式法则得
\begin{align}
D\overline{u}(x) = D\overline{u}'(\Phi(x)) D\Phi(x). \label{eq:ch05-chain-rule-uv}
\end{align}
对$\forall A\in \mathbb{R}^{n\times n}$,定义矩阵范数$\left\| A \right\| \triangleq \underset{v\in \mathbb{R} ^n\backslash \left\{ 0 \right\}}{\mathrm{sup}}\frac{\left\| Av \right\|}{\left\| v \right\|}$.下面证明 \(\|D\Phi(x)\|\) 在 \(\overline{W}\) 上关于 \(x\) 有界。因为 \(\Phi \in C^1\)，所以 \(D\Phi(x)\) 的每个元素 \(\frac{\partial \Phi_i}{\partial x_j}(x)\) 在 \(\overline{W}\) 上连续，从而矩阵值函数 \(x \mapsto D\Phi(x)\) 连续。
又由Cauchy-Schwarz不等式，对任意单位向量 \(v \in \mathbb{R}^n\)，有
\begin{align*}
\|Av\|^2 &= \sum_{i=1}^n \left( \sum_{j=1}^n a_{ij} v_j \right)^2 \leqslant \sum_{i=1}^n \left( \sum_{j=1}^n a_{ij}^2 \right) \left( \sum_{j=1}^n v_j^2 \right) = \sum_{i,j=1}^n a_{ij}^2.
\end{align*}
因此
\begin{align*}
\|A\| &\leqslant \left( \sum_{i,j=1}^n a_{ij}^2 \right)^{\frac{1}{2}}.
\end{align*}
于是当 \(x_k \to x\) 时，结合\cref{Functional Analysis-theorem:有界线性算子全体构成的空间完备}有
\begin{align*}
\bigl| \|D\Phi(x_k)\| - \|D\Phi(x)\| \bigr| \leqslant \|D\Phi(x_k) - D\Phi(x)\| &\leqslant \left( \sum_{i,j=1}^n \left( \frac{\partial \Phi_i}{\partial x_j}(x_k) - \frac{\partial \Phi_i}{\partial x_j}(x) \right)^2 \right)^{\frac{1}{2}} \to 0.
\end{align*}
故 \(x \mapsto \|D\Phi(x)\|\) 在 \(\overline{W}\) 上连续。由于 \(\overline{W}\) 是紧集，连续函数在紧集上必有最大值，记
\begin{align*}
C_1 &\triangleq \max_{x \in \overline{W}} \|D\Phi(x)\|,
\end{align*}
于是对一切 \(x \in \overline{W}\)，有 \(\|D\Phi(x)\| \leqslant C_1\)。由\cref{Functional Analysis-proposition:有界线性算子的基本不等式}，对任意 \(v \in \mathbb{R}^n\)，有
\begin{align*}
|v D\Phi(x)| \leqslant \|D\Phi(x)\| |v|.
\end{align*}
取 \(v = D\overline{u}'(\Phi(x))\)，并结合 \eqref{eq:ch05-chain-rule-uv}，得到
\begin{align}\label{eq::jioJFIOWJIOIJOjiojfwio-1}
|D\overline{u}(x)| = |D\overline{u}'(\Phi(x)) D\Phi(x)| \leqslant \|D\Phi(x)\| |D\overline{u}'(\Phi(x))| \leqslant C_1 |D\overline{u}'(\Phi(x))|.
\end{align}
同理可得,存在$C_2>0$,使得
\begin{align}\label{eq::jioJFIOWJIOIJOjiojfwio-2}
|Du'(x)| = |Du'(\Psi(x)) D\Psi(x)| \leqslant \|D\Psi(x)\| |Du'(\Psi(x))| \leqslant C_2 |Du'(\Psi(x))|.
\end{align}
利用\eqref{eq::jioJFIOWJIOIJOjiojfwio-1}\eqref{eq::jioJFIOWJIOIJOjiojfwio-2}式,再由变量替换 \(y = \Phi(x)\), \(x = \Psi(y)\) 及\eqref{eq::jifjiojiojiojiomfijsfiojjsifjioejiojmwfiejfiwmfiowijijwfio-1}\eqref{eq::jifjiojiojiojiomfijsfiojjsifjioejiojmwfiejfiwmfiowijijwfio-2}\eqref{eq::jifjiojiojiojiomfijsfiojjsifjioejiojmwfiejfiwmfiowijijwfio-3}式, 得
\begin{gather}
\int_W |D\overline{u}(x)|^p \mathrm{d}x \leqslant C_1^p \int_W |D\overline{u}'(\Phi(x))|^p \mathrm{d}x = C_1^p \int_B |D\overline{u}'(y)|^p |\det D\Psi(y)| \mathrm{d}y=C_1^p \int_B |D\overline{u}'(y)|^p \mathrm{d}y,\label{eq::jijioqwio93283r3r23-1}
\\
\int_{B^+}{\left| Du' \left( x \right) \right|^p\mathrm{d}x}\leqslant C_{2}^{p}\int_{B^+}{\left| Du' \left( \Psi \left( x \right) \right) \right|^p\mathrm{d}x}=C_{2}^{p}\int_{W\cap U}{\left| Du \left( y \right) \right|^p |\det D\Phi(y)|\mathrm{d}y}=C_{2}^{p}\int_{W\cap U}{\left| Du \left( y \right) \right|^p \mathrm{d}y}.\label{eq::jijioqwio93283r3r23-2}
\end{gather}
结合\eqref{eq::jijioqwio93283r3r23}\eqref{eq::jijioqwio93283r3r23-111}\eqref{eq::jijioqwio93283r3r23-1}\eqref{eq::jijioqwio93283r3r23-2}式可得
\begin{align}
\left\| \overline{u} \right\| _{W^{1,p}\left( W \right)}&=\left( \int_W{\left| \overline{u}\left( x \right) \right|^p\mathrm{d}x}+\int_W{\left| D\overline{u}\left( x \right) \right|^p\mathrm{d}x} \right) ^{\frac{1}{p}}\nonumber
\\
&\leqslant \left( \int_B{\left| \overline{u'}\left( y \right) \right|^p\mathrm{d}y}+C_{1}^{p}\int_B{\left| D\overline{u'}\left( y \right) \right|^p\mathrm{d}y} \right) ^{\frac{1}{p}}\nonumber\\
&\leqslant C_3\left\| \overline{u'} \right\| _{W^{1,p}\left( B \right)},
\label{eq:ch05-local-est-1}
\end{align}
\begin{align}
\left\| u^{\prime} \right\| _{W^{1,p}\left( B^+ \right)}&=\left( \int_{B^+}{\left| u^{\prime}\left( x \right) \right|^p\mathrm{d}x}+\int_{B^+}{\left| Du^{\prime}\left( x \right) \right|^p\mathrm{d}x} \right) ^{\frac{1}{p}}
\nonumber\\
&\leqslant \left( \int_{W\cap U}{\left| u\left( y \right) \right|^p\mathrm{d}y}+C_{2}^{p}\int_{W\cap U}{\left| Du\left( y \right) \right|^p\mathrm{d}y} \right) ^{\frac{1}{p}}
\nonumber\\
&\leqslant C_3\left\| u \right\| _{W^{1,p}\left( W\cap U \right)}\leqslant C_3\left\| u \right\| _{W^{1,p}\left( U \right)},\label{eq:ch05-local-est-2}
\end{align}
其中$C_3=\max \{ 1+C_{1},1+C_{2}\}$.再记$C'\triangleq CC_3^2$,将 \eqref{eq:ch05-ext-in-ball-2}\eqref{eq:ch05-local-est-1} \eqref{eq:ch05-local-est-2}式合并, 得到
\begin{align}
\left\| \overline{u} \right\| _{W^{1,p}\left( W \right)}\leqslant C_3\left\| \overline{u^{\prime}} \right\| _{W^{1,p}\left( B \right)}\leqslant CC_3\left\| u^{\prime} \right\| _{W^{1,p}\left( B^+ \right)}\leqslant CC_{3}^{2}\left\| u \right\| _{W^{1,p}\left( U \right)}=C^{\prime}\left\| u \right\| _{W^{1,p}\left( U \right)}.\label{eq:ch05-local-est-final}
\end{align}

若 \(p=\infty\)，先说明 \(L^\infty\) 范数在 \(C^1\) 微分同胚下不变。由\cref{lemma:正交平移变换为等距双射}和\cref{Real Analysis-lemma:Lipschitz映射保持零测集}可知，\(\Phi\) 与 \(\Psi\) 均把零测集映为零测集.设 \(f:B\to\mathbb{R}\),$\,g:W\cap U\to \mathbb{R}$均可测,下面证明
\begin{align}\label{eq::jiowiowijoijoijijojsiosioios3829237489}
\|f\circ\Phi\|_{L^\infty(W)}=\|f\|_{L^\infty(B)},\qquad
\|g\circ\Psi\|_{L^\infty(B^+)}=\|g\|_{L^\infty(W\cap U)}.
\end{align}
若 \(\|f\|_{L^\infty(B)}=\infty\)，则对任意 \(M>0\)，集合
\begin{align*}
N_M\triangleq\{y\in B:|f(y)|>M\}
\end{align*}
不是零测集。由 \(\Phi\) 把零测集映为零测集，故若 \(\Psi(N_M)\) 为零测集，则
\begin{align*}
N_M=\Phi(\Psi(N_M))
\end{align*}
为零测集，矛盾！因此 \(\Psi(N_M)\) 不是零测集。而
\begin{align*}
\{x\in W:|f(\Phi(x))|>M\}=\Phi^{-1}(N_M)=\Psi(N_M),
\end{align*}
故对任意 \(M>0\)，该集合不是零测集，所以
\begin{align*}
\|f\circ\Phi\|_{L^\infty(W)}=\infty.
\end{align*}
等号成立。

若 \(\|f\|_{L^\infty(B)}<\infty\)，记$M_0\triangleq\|f\|_{L^\infty(B)},$
则 \(|f|\leqslant M_0\) 几乎处处于 \(B\)。令
\begin{align*}
N\triangleq\{y\in B:|f(y)|>M_0\},
\end{align*}
则 \(N\) 为零测集。由 \(\Psi\) 把零测集映为零测集知\(\Psi(N)\) 为零测集。而
\begin{align*}
\{x\in W:|f(\Phi(x))|>M_0\}=\Phi^{-1}(N)=\Psi(N),
\end{align*}
故 \(|f\circ\Phi|\leqslant M_0\) 几乎处处于 \(W\)。因此
\begin{align}\label{eq::jiojfwiojiojwio3i2900384}
\|f\circ\Phi\|_{L^\infty(W)}\leqslant M_0=\|f\|_{L^\infty(B)}.
\end{align}
反过来，记$M_1\triangleq\|f\circ\Phi\|_{L^\infty(W)},$
则 \(|f\circ\Phi|\leqslant M_1\) 几乎处处于 \(W\)。令
\begin{align*}
N'\triangleq\{x\in W:|f(\Phi(x))|>M_1\}.
\end{align*}
则 \(N'\) 为零测集。由 \(\Phi\) 把零测集映为零测集知\(\Phi(N')\) 为零测集。而
\begin{align*}
\Phi(N')=\{y\in B:|f(y)|>M_1\},
\end{align*}
故 \(|f|\leqslant M_1\) 几乎处处于 \(B\)。因此
\begin{align}\label{eq::jiojfwiojiojwio3i2900382}
\|f\|_{L^\infty(B)}\leqslant M_1=\|f\circ\Phi\|_{L^\infty(W)}.
\end{align}
合并\eqref{eq::jiojfwiojiojwio3i2900384}\eqref{eq::jiojfwiojiojwio3i2900382}式，得
\begin{align*}
\|f\circ\Phi\|_{L^\infty(W)}=\|f\|_{L^\infty(B)}.
\end{align*}
把 \(f\) 换成 \(g\)，把 \(\Phi\) 换成 \(\Psi\)，并将 \(W\)换成$B^+$,\(B\)换成$W\cap U$,同理可得
\begin{align*}
\|g\circ\Psi\|_{L^\infty(B^+)}=\|g\|_{L^\infty(W\cap U)}.
\end{align*}

现在一方面,由 \(\overline u(x)=\overline u'(\Phi(x)),x\in W\)和\eqref{eq::jiowiowijoijoijijojsiosioios3829237489}式，得到
\begin{align*}
\|\overline u\|_{L^\infty(W)}
=\|\overline u'\circ\Phi\|_{L^\infty(W)}
=\|\overline u'\|_{L^\infty(B)}.
\end{align*}
再结合\eqref{eq::jioJFIOWJIOIJOjiojfwio-1}\eqref{eq::jiowiowijoijoijijojsiosioios3829237489}式可得
\begin{align*}
\|D\overline u\|_{L^\infty(W)}
\leqslant C_1\|D\overline u'\circ\Phi\|_{L^\infty(W)}
=C_1\|D\overline u'\|_{L^\infty(B)}.
\end{align*}
于是
\begin{align}\label{eq::jiwjiofijoswjiowiow2390jw89023890}
\left\| \overline{u} \right\| _{W^{1,\infty}\left( W \right)}=\left\| \overline{u} \right\| _{L^{\infty}\left( W \right)}+\left\| D\overline{u} \right\| _{L^{\infty}\left( W \right)}\leqslant \left\| \overline{u}^{\prime} \right\| _{L^{\infty}\left( B \right)}+C_1\left\| D\overline{u}^{\prime} \right\| _{L^{\infty}\left( B \right)}\leqslant C_3\left\| \overline{u}^{\prime} \right\| _{W^{1,\infty}\left( B \right)}.
\end{align}

另一方面，由 \(u'(y)=u(\Psi(y))\)，\(y\in B^+\)和\eqref{eq::jiowiowijoijoijijojsiosioios3829237489}式有
\begin{align}
\left\| u' \right\| _{L^{\infty}\left( B^+ \right)}=\left\| u\circ \Psi \right\| _{L^{\infty}\left( B^+ \right)}=\left\| u \right\| _{L^{\infty}\left( W\cap U \right)}
\end{align}
再结合\eqref{eq::jioJFIOWJIOIJOjiojfwio-2}\eqref{eq::jiowiowijoijoijijojsiosioios3829237489}式可得
\begin{align*}
\|Du'\|_{L^\infty(B^+)}
\leqslant C_2\|Du\circ\Psi\|_{L^\infty(B^+)}
=C_2\|Du\|_{L^\infty(W\cap U)}.
\end{align*}
从而
\begin{align}
\left\| u^{\prime} \right\| _{W^{1,\infty}\left( B^+ \right)}&=\left\| u^{\prime} \right\| _{L^{\infty}\left( B^+ \right)}+\left\| Du^{\prime} \right\| _{L^{\infty}\left( B^+ \right)}\leqslant \left\| u \right\| _{L^{\infty}\left( W\cap U \right)}+C_2\left\| Du \right\| _{L^{\infty}\left( W\cap U \right)}
\leqslant C_3\left\| u \right\| _{W^{1,\infty}\left( W\cap U \right)}.\label{eq::jiwjiofijoswjiowiow2390jw89023890-1}
\end{align}
将\eqref{eq:ch05-ext-in-ball-2}\eqref{eq::jiwjiofijoswjiowiow2390jw89023890}\eqref{eq::jiwjiofijoswjiowiow2390jw89023890-1}式合并，得到
\begin{align*}
\|\overline u\|_{W^{1,\infty}(W)}
\leqslant C_3\|\overline u'\|_{W^{1,\infty}(B)}
\leqslant C C_3\|u'\|_{W^{1,\infty}(B^+)}
\leqslant C C_3^2\|u\|_{W^{1,\infty}(W\cap U)}\leqslant C C_3^2\|u\|_{W^{1,\infty}(U)}=C'\|u\|_{W^{1,\infty}(U)}.
\end{align*}

综上，对任意给定的 \(u\in C^1(\overline{U})\) 以及任意 \(x_0\in\partial U\)，存在开集\(W\)，满足 \(x_0\in W\subset\subset V\)，并存在函数 \(\overline{u}\in C^1(W)\) 和只依赖于 \(p,U,V\) 的常数 \(C'>0\)，使得
\begin{gather}
\overline{u}=u\quad \text{a.e.}\, x\in W\cap U,\nonumber \\
\|\overline{u}\|_{W^{1,p}(W)}
\leqslant C'\|u\|_{W^{1,p}(U)},\quad 1\leqslant p\leqslant \infty.\label{eq::eEREFFwww}
\end{gather}

\noindent$\mathbf{Step}\bm{6}$.{\heiti 有限覆盖、单位分解与整体估计.}

由$\partial U$ 紧和$\mathbf{Step}\bm{5}$知，存在有限个点 $x_{i}^{0}\in\partial U$和包含$x_i^0$的开集 $W_{i}\subset \subset V$ 以及$u$ 到 $W_{i}$ 的延拓 $\overline{u}_{i}$（$i=1,\cdots,N$），使 $\partial U\subset\bigcup_{i=1}^{N}W_{i}$,且
\begin{align}
\|\overline{u}_i\|_{W^{1,p}(W_i)}
\leqslant C'\|u_i\|_{W^{1,p}(U)},\quad i=1,2,\cdots,N,\,1\leqslant p\leqslant \infty.\label{eq::eEREFFwww1}
\end{align}
因为 \(\partial U\) 紧且 \(\partial U\subseteq\bigcup_{i=1}^N W_i\)，其中每个 \(W_i\) 是开集。考虑
\begin{align}\label{eq:ioujiwjioKKKI}
K\triangleq U\setminus\bigcup_{i=1}^N W_i.
\end{align}
先证明 \(\overline{K}\subseteq U\)。若存在 \(x\in\overline{K}\cap\partial U\)，则由 \(\partial U\subseteq\bigcup_{i=1}^N W_i\) 知存在 \(j\in\{1,\cdots,N\}\) 使得 \(x\in W_j\)。因为 \(W_j\) 是开集，所以存在 \(\delta>0\) 使得 \(B(x,\delta)\subseteq W_j\)。又 \(x\in\overline{K}\)，故 \(B(x,\delta)\cap K\neq\varnothing\)。这与$K\subseteq U\setminus\bigcup_{i=1}^N W_i\subseteq \mathbb{R}^n\setminus W_j$矛盾!
因此 \(\overline{K}\cap\partial U=\varnothing\)。又由\eqref{eq:ioujiwjioKKKI}式和\subcref{Real Analysis-theorem:边界和开集基本性质}{Real Analysis-theorem:边界和开集基本性质-5}知\(\overline{K}\subseteq \overline{U}=U\cup \partial U\),故$\overline{K}\subseteq U$.又 \(\overline{K}\) 是闭集且 \(\overline{K}\subseteq\overline{U}\)，而 \(\overline{U}\) 紧，故 \(\overline{K}\) 紧。

若 \(K=\varnothing\)，则 \(U\subseteq\bigcup_{i=1}^{N}W_i\)。若 \(U\neq\varnothing\)，任取 \(x_*\in U\) 及 \(r>0\)，使得
\[
B(x_*,r)\subset\subset U,
\]
令
\[
W_0\triangleq B(x_*,r).
\]
则 \(W_0\subset\subset U\)，且
\[
U\subseteq W_0\cup\bigcup_{i=1}^{N}W_i.
\]
若 \(U=\varnothing\)，定理平凡，取 \(E=0\) 即可。

以下设 \(K\neq\varnothing\)。
因为 \(\overline{K}\) 紧且 \(\overline{K}\subseteq U\)，而 \(U\) 是有界开集,从而\(\overline{K}\)紧、$\mathbb{R}^n\setminus U$闭,两者不交,所以由\subcref{Real Analysis-proposition:点集间的距离的性质}{Real Analysis-proposition:点集间的距离的性质-3}知$\mathrm{dist}(\overline{K},\mathbb{R}^n\setminus U)>0$.令$\eta\triangleq\frac{1}{2}\mathrm{dist}(\overline{K},\mathbb{R}^n\setminus U)>0$,并定义
\begin{align*}
W_0\triangleq\{x\in\mathbb{R}^n:\mathrm{dist}(x,\overline{K})<\eta\}.
\end{align*}
则由距离函数连续知\(W_0\) 是开集，且
\begin{align*}
K\subseteq \overline{K}\subseteq \overline{W_0}\subseteq\{x\in\mathbb{R}^n:\mathrm{dist}(x,\overline{K})\leqslant\eta\}\subseteq U.
\end{align*}
因此 \(W_0\subset\subset U\)。

最后验证 \(\overline{U}\subseteq W_0\cup\bigcup_{i=1}^N W_i\)。任取 \(x\in \overline{U}\)。若 \(x\in\bigcup_{i=1}^N W_i\)，则已属于右边。若 \(x\notin\bigcup_{i=1}^N W_i\)，则 \(x\in K\subseteq\overline{K}\subseteq W_0\)。故无论如何都有 \(x\in W_0\cup\bigcup_{i=1}^N W_i\)。因此
\begin{align*}
\overline{U}\subseteq W_0\cup\bigcup_{i=1}^N W_i=\bigcup_{i=0}^N W_i.
\end{align*}
由\hyperref[corollary:紧集上的光滑单位分解定理]{紧集上的光滑单位分解定理}知,
对$i=0,1,\cdots,N$,存在 $\zeta_i\in C_c^\infty(W_i)$, 使得
\begin{align*}
0\leqslant\zeta_i\leqslant1,\quad
\sum\limits_{i=0}^{N}\zeta_i(x)=1,\quad\forall x\in \overline{U}.
\end{align*}
记$\overline{u}\triangleq\sum_{i=0}^{N}\zeta_{i}\overline{u}_{i}$，其中 $\overline{u}_{0}=u$。
因为 \(\zeta_i\in C_c^\infty(W_i)\)，所以 \(\zeta_i\) 具有紧支集。设
\[
K_i\triangleq \operatorname{spt}\zeta_i\subset\subset W_i.
\]
由于 \(\zeta_i\) 是连续函数，而 \(K_i\) 是紧集，所以 \(\zeta_i\) 在 \(K_i\) 上有界，即
\[
\|\zeta_i\|_{L^\infty(W_i)}=\sup_{x\in W_i}|\zeta_i(x)|=\max_{x\in K_i}|\zeta_i(x)|<\infty.
\]
同理，对每个 \(j=1,\cdots,n\)，偏导数 \(D_j\zeta_i\) 也是连续的。由于 \(\zeta_i\) 在 \(W_i\setminus K_i\) 上恒为零，所以 \(D_j\zeta_i\) 也在 \(W_i\setminus K_i\) 上恒为零，因此
\[
\operatorname{spt}D_j\zeta_i\subseteq K_i.
\]
于是 \(D_j\zeta_i\) 也在紧集 \(K_i\) 上连续，从而有界，即
\[
\|D_j\zeta_i\|_{L^\infty(W_i)}=\max_{x\in K_i}|D_j\zeta_i(x)|<\infty.
\]
因此
\[
\|D\zeta_i\|_{L^\infty(W_i)}\triangleq \max_{1\leqslant j\leqslant n}\|D_j\zeta_i\|_{L^\infty(W_i)}<\infty.
\]
所以
\[
M_i\triangleq\|\zeta_i\|_{L^\infty(W_i)}+\|D\zeta_i\|_{L^\infty(W_i)}<\infty.
\]
对每个\(i=0,1,\cdots,N\)，由\hyperref[theorem:ch05-弱导数性质-4]{Leibniz公式}得
\begin{align*}
D(\zeta_i \overline{u}_i) = \overline{u}_i D\zeta_i + \zeta_i D\overline{u}_i.
\end{align*}
因此当 \(1\leqslant p<\infty\) 时，由\hyperref[Basis of Analytics-proposition:凸性不等式]{凸性不等式\ref{Basis of Analytics-proposition:凸性不等式-2}}有
\begin{align*}
\left\| \zeta _i\overline{u}_i \right\| _{W^{1,p}\left( W_i \right)}&=\left( \left\| \zeta _i\overline{u}_i \right\| _{L^p\left( W_i \right)}^{p}+\left\| \overline{u}_iD\zeta _i+\zeta _iD\overline{u}_i \right\| _{L^p\left( W_i \right)}^{p} \right) ^{\frac{1}{p}}
\\
&\leqslant \left\| \zeta _i\overline{u}_i \right\| _{L^p\left( W_i \right)}+\left\| \overline{u}_iD\zeta _i+\zeta _iD\overline{u}_i \right\| _{L^p\left( W_i \right)}
\\
&\leqslant \left\| \zeta _i\overline{u}_i \right\| _{L^p\left( W_i \right)}+2^{p-1}\left( \left\| \overline{u}_iD\zeta _i \right\| _{L^p\left( W_i \right)}+\left\| \zeta _iD\overline{u}_i \right\| _{L^p\left( W_i \right)} \right) 
\\
&\leqslant M_i\left\| \overline{u}_i \right\| _{L^p\left( W_i \right)}+2^{p-1}M_i\left\| \overline{u}_iD\zeta _i \right\| _{L^p\left( W_i \right)}+2^{p-1}M_i\left\| \zeta _iD\overline{u}_i \right\| _{L^p\left( W_i \right)}
\\
&\leqslant 2^{p-1}M_i\left\| \overline{u}_i \right\| _{W^{1,p}\left( W_i \right)}.
\end{align*}
从而
\begin{align*}
\parallel \overline{u}\parallel _{W^{1,p}(\mathbb{R} ^n)}&\leqslant \sum_{i=0}^N{\parallel \zeta _i\overline{u}_i}\parallel _{W^{1,p}(W_i)}\leqslant \sum_{i=0}^N{2^{p-1}M_i\parallel \overline{u}_i\parallel _{W^{1,p}(W_i)}}
\\
&\leqslant 2^{p-1}M_0\parallel u\parallel _{W^{1,p}(U)}+\sum_{i=1}^N{2^{p-1}M_iC'}\parallel u\parallel _{W^{1,p}(U)}.
\end{align*}
于是取$C'' \triangleq 2^{p-1}M_0+C' \sum_{i=1}^N{2^{p-1}M_i},$则
\begin{align*}
\|\overline{u}\|_{W^{1,p}(\mathbb{R}^n)}\leqslant C''\|u\|_{W^{1,p}(U)},
\end{align*}
其中常数 \(C''\) 只依赖于 \(U,V,p,n\) 以及单位分解的选取,而单位分解的选取只依赖于\(U,V,p,n\),故常数 \(C''\) 只依赖于 \(U,V,p,n\).

当$p=\infty$时,有
\begin{align*}
\left\| \zeta _i\overline{u}_i \right\| _{W^{1,\infty}\left( W_i \right)}&\leqslant \underset{x\in W_i}{\mathrm{ess}\,\mathrm{sup}}\left| \zeta _i\overline{u}_i \right|+\underset{x\in W_i}{\mathrm{ess}\,\mathrm{sup}}\left| \overline{u}_iD\zeta _i \right|+\underset{x\in W_i}{\mathrm{ess}\,\mathrm{sup}}\left| \zeta _iD\overline{u}_i \right|
\\
&\leqslant \left\| \zeta _i \right\| _{L^{\infty}\left( W_i \right)}\underset{x\in W_i}{\mathrm{ess}\,\mathrm{sup}}\left| \overline{u}_i \right|+\left\| \zeta _i \right\| _{L^{\infty}\left( W_i \right)}\underset{x\in W_i}{\mathrm{ess}\,\mathrm{sup}}\left| D\overline{u}_i \right|+\left\| D\zeta _i \right\| _{L^{\infty}\left( W_i \right)}\underset{x\in W_i}{\mathrm{ess}\,\mathrm{sup}}\left| \overline{u}_i \right|
\\
&\leqslant \left( \left\| \zeta _i \right\| _{L^{\infty}\left( W_i \right)}+\left\| D\zeta _i \right\| _{L^{\infty}\left( W_i \right)} \right) \underset{x\in W_i}{\mathrm{ess}\,\mathrm{sup}}\left| \overline{u}_i \right|+\left\| \zeta _i \right\| _{L^{\infty}\left( W_i \right)}\underset{x\in W_i}{\mathrm{ess}\,\mathrm{sup}}\left| D\overline{u}_i \right|
\\
&\leqslant \left( \left\| \zeta _i \right\| _{L^{\infty}\left( W_i \right)}+\left\| D\zeta _i \right\| _{L^{\infty}\left( W_i \right)} \right) \left( \underset{x\in W_i}{\mathrm{ess}\,\mathrm{sup}}\left| \overline{u}_i \right|+\underset{x\in W_i}{\mathrm{ess}\,\mathrm{sup}}\left| D\overline{u}_i \right| \right) 
\\
&\leqslant M_{i}\left\| \overline{u}_i \right\| _{W^{1,\infty}\left( W_i \right)}.
\end{align*}
又因为$W_i\subset \subset V$，所以
\begin{align*}
\mathrm{spt}\,\overline{u}\subseteq \bigcup_{i=0}^N{\mathrm{spt}\,\zeta _i}\subseteq \bigcup_{i=0}^N{W_i}\subset \subset V.
\end{align*}
从而
\begin{align*}
\left\| \overline{u} \right\| _{W^{1,\infty}\left( \mathbb{R} ^n \right)}\leqslant \sum_{i=0}^N{\left\| \zeta _i\overline{u}_i \right\| _{W^{1,\infty}\left( W_i \right)}}\leqslant \sum_{i=0}^N{M_{i}}\left\| \overline{u}_i \right\| _{W^{1,\infty}\left( W_i \right)}
\\
\leqslant M_{0}\left\| u \right\| _{W^{1,\infty}\left( U \right)}+\sum_{i=1}^N{M_{i}C\prime}\left\| u \right\| _{W^{1,\infty}\left( U \right)}.
\end{align*}
于是取$C'' \triangleq M_0 + C' \sum\limits_{i=1}^N M_i,$则
\begin{align*}
\|\overline{u}\|_{W^{1,p}(\mathbb{R}^n)}\leqslant C''\|u\|_{W^{1,p}(U)},
\end{align*}
其中常数 \(C''\) 只依赖于 \(U,V,p,n\) 以及单位分解的选取,而单位分解的选取只依赖于\(U,V,p,n\),故常数 \(C''\) 只依赖于 \(U,V,p,n\).综上,对$1\leqslant p\leqslant \infty$,都存在只依赖于 \(U,V,p,n\)的常数$C''>0$,使得
\begin{align}\label{eq::fjiwjio233--12}
\|\overline{u}\|_{W^{1,p}(\mathbb{R}^n)}\leqslant C''\|u\|_{W^{1,p}(U)},
\end{align}

\noindent$\mathbf{Step}\bm{7}$.{\heiti 由光滑逼近定义 \(E\)（\(1\leqslant p<\infty\)）.}

对于 \(u\in C^{\infty}(\overline{U})\)，由$\mathbf{Step 1-6}$构造的 \(\overline{u}\) 定义$E_0 u\triangleq \overline{u},$
则由$\mathbf{Step 1-6}$的构造都是线性的可知\(E_0:C^{\infty}(\overline{U})\to W^{1,p}(\mathbb{R}^n)\) 是线性的，且由\eqref{eq::fjiwjio233--12}式有
\begin{align}\label{eqJIOFJIOWJIO323u89rfu3}
\|E_0 u\|_{W^{1,p}(\mathbb{R}^n)}\leqslant C''\|u\|_{W^{1,p}(U)}.
\end{align}

现设 \(1\leqslant p<\infty\)，\(u\in W^{1,p}(U)\)。由\cref{theorem:ch05-直到边界光滑的函数整体逼近定理}，可选 \(u_m\in C^{\infty}(\overline{U})\) 在 \(W^{1,p}(U)\) 中收敛于 \(u\)。由\eqref{eqJIOFJIOWJIO323u89rfu3}式与 \(E_0\) 的线性性可得
\begin{align*}
\|E_0 u_m-E_0 u_l\|_{W^{1,p}(\mathbb{R}^n)}
&\leqslant C''\|u_m-u_l\|_{W^{1,p}(U)}.
\end{align*}
故 \(\{E_0 u_m\}_{m=1}^{\infty}\) 是 \(W^{1,p}(\mathbb{R}^n)\) 中的 Cauchy 列。由\cref{theorem:ch05-WkpBanach}，\(W^{1,p}(\mathbb{R}^n)\) 是 Banach 空间，故存在 \(\overline{u}\in W^{1,p}(\mathbb{R}^n)\) 使得
\begin{align*}
E_0 u_m\to \overline{u}\quad \text{于 } W^{1,p}(\mathbb{R}^n).
\end{align*}
定义$E u\triangleq \overline{u}.$
下证该$E$是良定义的。假设 \(\{v_m\}_{m=1}^{\infty}\subseteq C^{\infty}(\overline{U})\) 也在 \(W^{1,p}(U)\) 中收敛于 \(u\)。则对任意 \(m\)，
\begin{align*}
\|E_0 u_m-E_0 v_m\|_{W^{1,p}(\mathbb{R}^n)}
\leqslant C''\|u_m-v_m\|_{W^{1,p}(U)}.
\end{align*}
而
\begin{align*}
\|u_m-v_m\|_{W^{1,p}(U)}
\leqslant \|u_m-u\|_{W^{1,p}(U)}+\|u-v_m\|_{W^{1,p}(U)}\to 0\quad (m\to\infty).
\end{align*}
故
\begin{align*}
\|E_0 u_m-E_0 v_m\|_{W^{1,p}(\mathbb{R}^n)}\to 0\quad (m\to\infty).
\end{align*}
由于 \(\{E_0 u_m\}\) 与 \(\{E_0 v_m\}\) 分别收敛于 \(\overline{u}\) 与 \(\overline{v}\)，于是
\begin{align*}
\|\overline{u}-\overline{v}\|_{W^{1,p}(\mathbb{R}^n)}
=\lim\limits_{m\to\infty}\|E_0 u_m-E_0 v_m\|_{W^{1,p}(\mathbb{R}^n)}=0.
\end{align*}
所以 \(\overline{u}=\overline{v}\)。因此 \(E u\) 不依赖于逼近序列的选取。

由 \(E_0\) 的线性性和极限的线性性知\(E\) 是线性的,并且由\eqref{eqJIOFJIOWJIO323u89rfu3}式知
\begin{align*}
\|E u\|_{W^{1,p}(\mathbb{R}^n)}
&=\lim\limits_{m\to\infty}\|E_0 u_m\|_{W^{1,p}(\mathbb{R}^n)}
\leqslant C''\lim\limits_{m\to\infty}\|u_m\|_{W^{1,p}(U)}
=C''\|u\|_{W^{1,p}(U)}.
\end{align*}
故 \(E\) 有界。此外，再证明 \(Eu=u\) a.e. 于 \(U\)。因为 \(u_m\in C^\infty(\overline U)\) 且 \(u_m\to u\) 于 \(W^{1,p}(U)\)，所以 \(u_m\to u\) 于 \(L^p(U)\)。由 \(E_0\) 的定义，对每个 \(m\) 有 \(E_0 u_m=u_m\) a.e. 于 \(U\)，因此
\begin{align*}
E_0 u_m|_U=u_m\to u\quad \text{于 } L^p(U).
\end{align*}
另一方面，由 \(Eu\) 的定义，\(E_0 u_m\to Eu\) 于 \(W^{1,p}(\mathbb R^n)\)，从而限制在 \(U\) 上，有
\begin{align*}
E_0 u_m|_U\to Eu|_U\quad \text{于 } L^p(U).
\end{align*}
由 \(L^p(U)\) 中极限的唯一性，得到
\begin{align*}
Eu|_U=u\quad \text{于 } L^p(U),
\end{align*}
即 \(Eu=u\) a.e. 于 \(U\)。

再证 \(\operatorname{spt}Eu\subseteq V\)。由第6步的构造，所有 \(E_0 u_m\) 的支集都含于紧集
\begin{align*}
K\triangleq \bigcup_{i=0}^N \overline{W_i}.
\end{align*}
因为 \(W_i\subset\subset V\)，所以 \(K\subset\subset V\)。于是对每个 \(m\)，有
\begin{align*}
E_0 u_m=0\quad \text{a.e. 于 } \mathbb R^n\setminus K.
\end{align*}
由 \(E_0 u_m\to Eu\) 于 \(W^{1,p}(\mathbb R^n)\)，限制在 \(\mathbb R^n\setminus K\) 上得到
\begin{align*}
E_0 u_m|_{\mathbb R^n\setminus K}\to Eu|_{\mathbb R^n\setminus K}\quad \text{于 } L^p(\mathbb R^n\setminus K).
\end{align*}
但 \(E_0 u_m|_{\mathbb R^n\setminus K}=0\) 于 \(L^p(\mathbb R^n\setminus K)\)，故
\begin{align*}
Eu|_{\mathbb R^n\setminus K}=0\quad \text{于 } L^p(\mathbb R^n\setminus K),
\end{align*}
即 \(Eu=0\) a.e. 于 \(\mathbb R^n\setminus K\)。因此
\begin{align*}
\operatorname{spt}Eu\subseteq K\subset\subset V.
\end{align*}
综上,定理对 \(1\leqslant p<\infty\) 成立。

\noindent$\mathbf{Step}\bm{8}$.{\heiti \(p=\infty\) 的 Lipschitz 延拓.}

设 \(u\in W^{1,\infty}(U)\)。固定 \(x_0\in\partial U\)。取 $\mathbf{Step 4}$中的坐标邻域 \(W\subset\subset V\) 和 \(C^1\) 微分同胚 \(\Phi:W\to B\)，\(\Psi=\Phi^{-1}\)，使得
\begin{align*}
\Phi(W\cap U)=B^+,\qquad
\Phi(W\cap\partial U)=B\cap\{y_n=0\}.
\end{align*}
定义
\begin{align*}
u'(y)\triangleq u(\Psi(y)),\qquad y\in B^+.
\end{align*}
由 \(u\in W^{1,\infty}(U)\) 及 \(\Psi\) 为 \(C^1\) 微分同胚，由弱链式法则，有
\begin{align*}
u'\in W^{1,\infty}(B^+),
\end{align*}
并且存在常数 \(C>0\)，使得
\begin{align}
\|u'\|_{W^{1,\infty}(B^+)}
\leqslant C\|u\|_{W^{1,\infty}(W\cap U)}
\leqslant C\|u\|_{W^{1,\infty}(U)}.
\label{eq:u'-W1infty估计}
\end{align}
由于 \(B^+\) 是凸集，由\cref{theorem:ch05-W1infty局部刻画}知，\(u'\) 在 \(B^+\) 上 Lipschitz，且其 Lipschitz 常数不超过
\begin{align*}
L'\triangleq\|Du'\|_{L^\infty(B^+)}.
\end{align*}

按$\mathbf{Step 1}$的公式定义
\begin{align*}
\overline{u'}(y)\triangleq
\begin{cases}
u'(y), & y\in B^+,\\[4pt]
-3u'(y',-y_n)+4u'\left(y',-\dfrac{y_n}{2}\right), & y\in B^-.
\end{cases}
\end{align*}
在 \(y_n=0\) 上，两侧都等于 \(u'(y',0)\)，故$\overline{u'}$是良定义的。

下面证明 \(\overline{u'}\) 在 \(B\) 上 Lipschitz。在 \(B^+\) 上，\(\overline{u'}=u'\)，Lipschitz 常数为 \(L'\)。在 \(B^-\) 上，令
\begin{align*}
T_1(y)\triangleq(y',-y_n),\qquad
T_2(y)\triangleq\left(y',-\frac{y_n}{2}\right).
\end{align*}
\(T_1,T_2\) 都是 \(1\)-Lipschitz 的仿射映射，且
\begin{align*}
T_1(B^-)\subseteq B^+,\qquad
T_2(B^-)\subseteq B^+.
\end{align*}
因此对 \(y,z\in B^-\)，
\begin{align*}
|\overline{u'}(y)-\overline{u'}(z)|
&\leqslant3|u'(T_1y)-u'(T_1z)|
+4|u'(T_2y)-u'(T_2z)|\\
&\leqslant3L'|T_1y-T_1z|+4L'|T_2y-T_2z|
\leqslant7L'|y-z|.
\end{align*}
故 \(\overline{u'}\) 在 \(B^-\) 上 \(7L'\)-Lipschitz。

若 \(y\in B^+\)，\(z\in B^-\)，令 \(w\) 为线段 \(yz\) 与 \(\{y_n=0\}\) 的交点。由于 \(u'\) 在 \(B^+\) 上 Lipschitz，它在 \(B^0\) 上有连续延拓，且
\begin{align*}
\overline{u'}(w)=u'(w).
\end{align*}
于是
\begin{align*}
|\overline{u'}(y)-\overline{u'}(z)|
\leqslant|\overline{u'}(y)-u'(w)|+|u'(w)-\overline{u'}(z)|\leqslant L'|y-w|+7L'|w-z|
\leqslant7L'|y-z|.
\end{align*}
因此 \(\overline{u'}\) 在 \(B\) 上 \(7L'\)-Lipschitz。

由 \cref{theorem:ch05-W1infty局部刻画}，\(\overline{u'}\in W^{1,\infty}(B)\)。又由差商估计，有
\begin{align*}
\|D\overline{u'}\|_{L^\infty(B)}\leqslant7L'.
\end{align*}
同时
\begin{align*}
|\overline{u'}|\leqslant7\|u'\|_{L^\infty(B^+)}.
\end{align*}
于是
\begin{align}
\|\overline{u'}\|_{W^{1,\infty}(B)}
\leqslant C\|u'\|_{W^{1,\infty}(B^+)}.
\label{eq:延拓u'-W1infty}
\end{align}

令
\begin{align*}
\overline u(x)\triangleq\overline{u'}(\Phi(x)),\qquad x\in W.
\end{align*}
因为 \(\overline{u'}\) Lipschitz、\(\Phi\in C^1\) 且 \(D\Phi\) 在 \(\overline W\) 上有界，所以 \(\overline u\) 在 \(W\) 上 Lipschitz，从而 \(\overline u\in W^{1,\infty}(W)\)。由 Lipschitz 复合和差商估计，
\begin{align}
\|\overline u\|_{W^{1,\infty}(W)}
\leqslant C\|\overline{u'}\|_{W^{1,\infty}(B)}
\leqslant C\|u'\|_{W^{1,\infty}(B^+)}
\leqslant C\|u\|_{W^{1,\infty}(W\cap U)}
\leqslant C\|u\|_{W^{1,\infty}(U)}.
\label{eq:局部延拓-W1infty}
\end{align}
若 \(x\in W\cap U\)，则 \(\Phi(x)\in B^+\)，于是
\begin{align*}
\overline u(x)
=\overline{u'}(\Phi(x))
=u'(\Phi(x))
=u(\Psi(\Phi(x)))
=u(x).
\end{align*}
故
\begin{align*}
\overline u\big|_{W\cap U}=u.
\end{align*}

沿用$\mathbf{Step 6}$中的有限覆盖 \(W_0,\cdots ,W_N\subset\subset V\)，其中
\begin{align*}
\overline U\subseteq\bigcup_{i=0}^N W_i,
\end{align*}
并对每个 \(i\geqslant1\)，\(\overline u_i\) 为上述局部延拓；令 \(\overline u_0=u\)。由\cref{corollary:紧集上的光滑单位分解定理}，存在
\begin{align*}
\zeta_i\in C_c^\infty(W_i),\qquad
0\leqslant\zeta_i\leqslant1,\qquad
\sum_{i=0}^N\zeta_i(x)=1,\quad x\in\overline U.
\end{align*}
定义
\begin{align*}
\overline u\triangleq\sum_{i=0}^N\zeta_i\overline u_i.
\end{align*}
每个 \(\zeta_i\overline u_i\) 都紧支于 \(W_i\) 且属于 \(W^{1,\infty}(\mathbb R^n)\)，故 \(\overline u\in W^{1,\infty}(\mathbb R^n)\)。由 \eqref{eq:局部延拓-W1infty} 及 Leibniz 公式，
\begin{align}
\|\overline u\|_{W^{1,\infty}(\mathbb R^n)}
&\leqslant\sum_{i=0}^N\|\zeta_i\overline u_i\|_{W^{1,\infty}(W_i)}
\leqslant C\sum_{i=0}^N\|\overline u_i\|_{W^{1,\infty}(W_i)}
\leqslant C\|u\|_{W^{1,\infty}(U)},
\label{eq:整体延拓-W1infty}
\end{align}
其中最后一步必要时增大 \(C\)，使其吸收 \(N+1\) 及所有 \(\zeta_i\) 的常数。

又因为 \(W_i\subset\subset V\)，所以
\begin{align*}
\operatorname{spt}\overline u
\subseteq\bigcup_{i=0}^N\operatorname{spt}\zeta_i
\subseteq\bigcup_{i=0}^N W_i
\subset\subset V.
\end{align*}
并且由 \(\sum_i\zeta_i=1\) 于 \(\overline U\)、\(\overline u_i=u\) 于 \(W_i\cap U\)，得到
\begin{align*}
\overline u|_U=u.
\end{align*}

定义
\begin{align*}
Eu\triangleq\overline u.
\end{align*}
由于反射公式和单位分解都是线性的，\(u\mapsto Eu\) 线性；由 \eqref{eq:整体延拓-W1infty}式，\(E\) 有界。故 \(p=\infty\) 时定理成立。
\end{proof}

\begin{remark}\label{remark:ch05-延拓注}
\begin{enumerate}[(i)]
\item 现设 $\partial U$ 为 $C^{2}$。则上面构造的延拓算子 $E$ 也是从 $W^{2,p}(U)$ 到 $W^{2,p}(\mathbb{R}^{n})$ 的有界线性算子。为看到这一点，首先注意在证明的步骤 3、4 中，尽管 $\overline{u}$ 一般不是 $C^{2}$，但它确实属于 $W^{2,p}(B)$。我们还可以由定义\eqref{align:ch05-反射延拓}得到
\begin{align*}
\|\overline{u}\|_{W^{2,p}(B)}\leqslant C\|u\|_{W^{2,p}(B^{+})},
\end{align*}
如前所述，只要 $\partial U$ 为 $C^{2}$，我们由此导出估计
\begin{align}
\|Eu\|_{W^{2,p}(\mathbb{R}^{n})}\leqslant C\|u\|_{W^{2,p}(U)},
\end{align}
其中常数 $C$ 只依赖于 $U$、$V$、$n$ 与 $p$。我们稍后将用到这些观察。
\item 上述构造并未为 $k>2$ 的Sobolev空间 $W^{k,p}(U)$ 提供延拓。这需要更复杂的高阶反射技巧。
\end{enumerate}
\end{remark}



\end{document}